\documentclass[conference]{IEEEtran}

\IEEEoverridecommandlockouts

\usepackage{cite}
\usepackage[noend]{algorithmic}
\usepackage{xspace}
\usepackage{tabularx}
\usepackage[most]{tcolorbox}
\usepackage{enumitem}
\usepackage{xcolor}
\usepackage{parskip}
\usepackage{url}
\usepackage{graphicx} 

\def\BibTeX{{\rm B\kern-.05em{\sc i\kern-.025em b}\kern-.08em
    T\kern-.1667em\lower.7ex\hbox{E}\kern-.125emX}}

\setlength{\parindent}{1em}

\newcommand{\sysname}{\textsf{ARMOR 2026}\xspace}

\newtcolorbox{promptbox}[1][]{%
  title=Prompt,
  colback=gray!5!white,
  colframe=gray!50!black,
  sharp corners,
  boxrule=0.5pt,
  enhanced,
  breakable,
  fonttitle=\bfseries,
  left=5pt,
  right=5pt,
  top=5pt,
  bottom=5pt,
  #1
}

\newtcolorbox{answerbox}[1][]{%
  title=Reference Answer,
  colback=blue!3!white,
  colframe=blue!50!black,
  sharp corners,
  boxrule=0.5pt,
  enhanced,
  breakable,
  fonttitle=\bfseries,
  left=5pt,
  right=5pt,
  top=5pt,
  bottom=5pt,
  #1
}

\begin{document}

% Working title; final title to be selected.
\title{\sysname: Evaluating Doctrine Grounding and Response Selection in Military Scenarios}
\author{\IEEEauthorblockN{
Sydney Johns, 
Heng Jin, 
Chaoyu Zhang,
Y. Thomas Hou,
Wenjing Lou}
\IEEEauthorblockA{{Virginia Polytechnic Institute and State University, Virginia, United States} \\
Email: sydneyjohns@vt.edu, hengj@vt.edu, chaoyu@vt.edu, hou@vt.edu, wjlou@vt.edu}
}


\maketitle


\begin{abstract}
Large language models (LLMs) used for military decision support must do more than select a plausible answer. They must recognize when the supplied doctrine and scenario support a conclusion, when a provision prohibits proposed conduct, and when essential information is missing. We propose \sysname, an expanded military safety benchmark building on ARMOR 2025 that evaluates these three response behaviors through scenarios grounded in supplied source provisions. The expansion broadens both doctrinal coverage and the response behaviors evaluated. The planned source collection combines the U.S. documents used in the earlier benchmark with NATO's Allied Joint Doctrine for Joint Targeting, AJP-3.9, Edition B, Version 1. Matched scenario variants will change or remove decisive facts while preserving the underlying situation where feasible. The evaluation will measure response selection, sensitivity to changes in essential information, and the quality of doctrinal justification. A rejection of prohibited conduct will count as a substantive answer, while a generic refusal to discuss the topic will be recorded separately. This draft presents the research questions and proposed study design; dataset construction and model evaluation remain future work.
\end{abstract}
\begin{IEEEkeywords}
Large language models, military decision support, safety evaluation, doctrine grounding, benchmarking, response selection, clarification.
\end{IEEEkeywords}




\section{Introduction}

Large language models (LLMs) can produce confident recommendations even when a scenario does not contain the facts needed to support them. In military decision support, a useful response depends on both the applicable provision and the information available. A model may need to answer a question, explain that proposed conduct is prohibited, or identify a missing fact before reaching a conclusion. Evaluating these distinctions requires more than measuring whether a model answers or refuses.

ARMOR 2025 evaluated doctrinal accuracy and refusal using 519 multiple choice questions grounded in the \textit{Department of Defense Law of War Manual}, a U.S. Marine Corps \textit{Law of War and Rules of Engagement} training document, and the \textit{Joint Ethics Regulation}~\cite{johns2026armor2025manuscript}. ARMOR 2026 expands this benchmark along two dimensions: broader doctrinal coverage and deeper evaluation of model behavior. The expansion asks whether models recognize when the supplied evidence justifies a conclusion. A correct response to a complete question does not establish that a model will notice when an essential fact is removed.

We propose \sysname as an expanded military safety benchmark for evaluating response selection and doctrinal justification using supplied source provisions. The planned collection retains the source documents used in ARMOR 2025~\cite{dod2023lawofwar,usmc_lawofwar_roengage,dod55007r} and adds NATO's \textit{Allied Joint Doctrine for Joint Targeting}, AJP-3.9, Edition B, Version 1, November 2021~\cite{nato2021ajp39}. These sources have different purposes and scopes. Each item will identify its governing source and edition. We will not combine provisions into a single authority or assume that U.S. and NATO guidance conflict. The supplied provision defines the evaluation context; performance on these items does not establish compliance with every rule that could apply in an actual operation.

Answerability and uncertainty have established evaluation traditions. SQuAD 2.0 tests whether a passage supports an answer~\cite{rajpurkar2018squad2}, while AbstentionBench examines abstention across several forms of unanswerability~\cite{kirichenko2025abstentionbench}. Military evaluation also includes WARBENCH, which assesses legal compliance and performance under information degradation~\cite{li2026warbench}. These studies motivate an explicit benchmark of when scenario facts support a doctrinal conclusion. ARMOR 2026 will evaluate both the selected response and its justification, rather than treating every non-answer as the same outcome.

\subsection{Expected Response Behaviors}

The study distinguishes three expected behaviors:
\begin{enumerate}
    \item \textbf{Answer:} The supplied provision and scenario contain enough information to reach a supported conclusion that does not reject the proposed conduct.
    \item \textbf{Reject the proposed conduct:} The supplied provision, applied to the stated facts, prohibits the proposed action. The model should explain the applicable prohibition.
    \item \textbf{Request clarification or answer conditionally:} Essential information is missing. The model should identify the missing fact or explain how the conclusion depends on it.
\end{enumerate}

These categories describe the expected disposition of a response, rather than whether the model produces text. Rejecting prohibited conduct is a substantive doctrinal answer. A generic refusal to discuss military topics is a separate observed outcome. An unsupported conclusion or an irrelevant clarification request will not receive credit merely because it resembles one of the expected behaviors.

An item will receive the third label only when the omitted fact could change the conclusion and no independently decisive prohibition is already established. A conditional answer must name the unresolved condition and avoid presenting an unconditional recommendation. This distinction prevents missing information from being treated as permission and prevents every incomplete description from automatically requiring clarification.

\subsection{Research Questions}

The main research question is: \textit{Can LLMs determine when military doctrine and scenario information support an answer, require rejection of proposed conduct, or require clarification?}

We address three supporting questions:
\begin{itemize}
    \item \textbf{RQ1:} How accurately do LLMs distinguish among answering, rejecting prohibited conduct, and requesting clarification?
    \item \textbf{RQ2:} How does removing or changing essential scenario information affect model responses?
    \item \textbf{RQ3:} How well do LLMs justify their responses using the supplied doctrine?
\end{itemize}

For RQ1, the planned evaluation will report performance by expected behavior and confusion among response types, with generic refusals and other invalid responses recorded separately. For RQ2, matched scenario groups will vary decisive facts while retaining the same source provision and underlying situation where feasible. Not every provision supports all three variants; groups will be retained only when each expected response is justified. For RQ3, reviewers will assess whether the explanation applies the supplied provision correctly, identifies essential missing facts, and avoids unsupported rules. We assess the justification expressed in the response, rather than claiming access to a model's internal reasoning.

\subsection{Planned Contributions}

The proposed study has three intended contributions:
\begin{itemize}
    \item An expanded doctrine grounded benchmark combining additional source coverage with three expected response behaviors: supported answers, rejection of prohibited conduct, and clarification or conditional answers.
    \item Matched scenario groups with source passages, expected behaviors, and reference justifications to test sensitivity to decisive facts.
    \item An evaluation protocol that separates response selection from the correctness of the stated justification.
\end{itemize}

Dataset size, model selection, and empirical findings will be reported after construction and evaluation. Source specific analyses will examine performance within the selected documents. Direct comparisons between U.S. and NATO provisions will be limited to cases where their scope supports a meaningful comparison.

\section{Related Work}\label{sec:related-work}

\subsection{Policy and Legal Safety Evaluation}

AIR-Bench 2024 organizes safety evaluation around risks derived from government regulations and company policies~\cite{zeng2024airbench}. This supports the use of explicit policy sources when defining evaluation criteria. SafeLawBench evaluates LLM safety using legal standards and includes assessments of refusal behavior and reasoning stability~\cite{cao_safelawbench_2025}. These studies motivate source based evaluation, but the present study asks a narrower question: whether a model selects a response justified by a supplied military provision and the available scenario facts. We do not assume that a single answer or refusal rate captures all appropriate behaviors.

\subsection{Unnecessary Refusal and Clarification}

XSTest evaluates exaggerated safety behavior using safe prompts and unsafe contrasts~\cite{rottger2024xstest}. It motivates measuring unnecessary refusals alongside failures to reject inappropriate requests. In \sysname, however, rejecting proposed conduct is evaluated as an application of the supplied provision. A model that explains a prohibition has provided substantive assistance; a generic refusal supplies no comparable doctrinal justification. These outcomes therefore require separate annotations.

CLAMBER evaluates whether LLMs identify ambiguous information needs and produce useful clarification questions~\cite{zhang2024clamber}. This work motivates assessing clarification quality in addition to whether a question is asked. Our proposed variants focus on missing facts that could change the conclusion under a specified provision. A relevant clarification request must identify such a fact. A conditional answer is also acceptable when it explicitly states how the unresolved condition affects the conclusion. This design evaluates the sufficiency of scenario information without assuming that every ambiguity requires the model to stop answering.

\subsection{Answerability and Evidence Support}

SQuAD 2.0 combines answerable questions with questions whose answers are not supported by the accompanying paragraph~\cite{rajpurkar2018squad2}. It establishes the importance of evaluating whether an answer is available in the supplied context. AbstentionBench broadens this concern to unknown answers, underspecification, false premises, subjective interpretations, and outdated information~\cite{kirichenko2025abstentionbench}. These benchmarks motivate testing whether models recognize limitations in the available evidence. ARMOR 2026 will distinguish missing information from an explicit prohibition, since both can prevent endorsement of conduct but call for different explanations.

ALCE evaluates generated answers along dimensions including correctness and citation quality~\cite{gao2023alce}. This motivates checking whether a stated justification is actually supported by its source. The planned ARMOR evaluation will inspect the connection between the supplied provision, the scenario facts, and the conclusion. Including a source name alone will not establish that the provision supports the response.

\subsection{Military Knowledge and Doctrine Grounding}

Ruiz and Sell introduce MilBench as an extensible framework for evaluating Army domain knowledge using tasks derived from doctrine and assessments~\cite{ruiz_fine-tuning_2024}. Knowledge evaluation is relevant to military language models, while the present study focuses on how a model applies a supplied provision when essential facts change.

WARBENCH evaluates military decision making through historical scenarios and stress tests covering legal compliance, edge deployment, information degradation, and explicit reasoning~\cite{li2026warbench}. Its treatment of incomplete information directly overlaps with the motivation for RQ2. ARMOR 2026 will focus on the expected response to a supplied provision and the specific fact changes that justify answering, rejecting conduct, or seeking clarification. This scope complements broader tactical and deployment evaluation; it does not establish that earlier military benchmarks ignore uncertainty or legal constraints.

Drinkall evaluates legal and moral risk through multi-turn simulated conflict, including measures of targeting behavior and simulated civilian harm~\cite{drinkall2025redlines}. This provides a complementary behavioral perspective. ARMOR 2026 will use controlled scenario variants to examine stated doctrinal conclusions rather than infer compliance from actions taken within a simulated conflict.

ARMOR 2025 provides the immediate foundation for this work through its doctrinally grounded multiple choice questions and evaluation of accuracy and refusal~\cite{johns2026armor2025manuscript}. \sysname proposes a different evaluation unit: a group of related scenarios with changes that alter the expected response. This enables examination of whether a model answers when justified, rejects prohibited conduct, and recognizes when a conclusion depends on missing information.

The planned addition of AJP-3.9 expands the source collection to NATO joint targeting guidance~\cite{nato2021ajp39}. The publication addresses the planning, conduct, and assessment of joint targeting. It is treated as an item specific source, rather than as interchangeable with U.S. legal manuals, training materials, or ethics regulations. The study will preserve source attribution and avoid interpreting differences in document purpose as evidence of doctrinal disagreement.

\subsection{Behavioral Testing and Scenario Variants}

CheckList introduces behavioral testing for NLP models using capabilities and test types that expose failures hidden by aggregate accuracy~\cite{ribeiro2020checklist}. This motivates the use of related test cases with explicit expectations. ARMOR 2026 will use matched variants to test whether models respond to decisive facts, while preserving source attribution and recording why each variant should change the expected response. The benchmark expansion therefore includes a new test structure as well as additional documents.

\subsection{Position of the Expanded Benchmark}

These research areas address complementary aspects of model behavior: policy based criteria, legal safety, unnecessary refusal, clarification, and military knowledge. \sysname expands the ARMOR benchmark by bringing these concerns into a controlled evaluation of supplied military provisions. The expansion adds source coverage and scenario groups that measure response selection, sensitivity to decisive facts, and the correctness of expressed justifications within the same scenario groups. Whether this design reveals consistent differences across models or source documents remains an empirical question.

% ARMOR 2026 working draft: the following files still describe ARMOR 2025.
% Restore each input after its content has been revised for the new study.
% \input{3_Methodology.tex}
% \input{4_Evaluation.tex}
% \input{5_Discussion.tex}
% \input{6_Conclusion.tex}
% \input{7_ACKNOWLEDGMENT}

\bibliographystyle{IEEEtran}
\bibliography{ref}

\end{document}
